% The look of the book: closed-source's paper (black and red) with picasso's one-line drawing on
% the title page. Only packages of picasso's `paper` lab are used, and fonts of its TeX tree.
% Taken in part from marola-dev/marola-research's format/header.tex at 51fae22 (MIT, same author).
\usepackage{graphicx}
\usepackage{xcolor}
% The equations of the title page: ≤, π and ² print only in a font that has them; DejaVu Sans Mono,
% by file name, from the TeX tree of the pinned toolchain.
\newfontfamily\eqfont{DejaVuSansMono.ttf}
\definecolor{bookred}{HTML}{A3201C}
\definecolor{bookink}{HTML}{111111}
\color{bookink}
% The title page: two plots made of equations, side by side, each with its equations as its
% caption (build/plots/cover.tex, written by `just plots`), in place of a drawing.
\makeatletter
\let\nash@maketitle\@maketitle
\def\@maketitle{%
  \nash@maketitle
  \IfFileExists{cover.tex}{\noindent\input{cover.tex}}{}
  \vskip 0.6em
  \begin{center}\small\color{bookred}an open book in progress. the voice is ``Nash'', a persona of the author;\\
  no sentence in it is the mathematician John Forbes Nash Jr.'s.\end{center}
  \vskip 1em}
\makeatother
